\examplegloss{fibre-product-of-stacks}{Fibre product of stacks}{
\begin{enumerate}[label=(\arabic*),leftmargin=*,itemsep=2pt,topsep=3pt]
\item For ordinary schemes, $\A^1\times_{\Spec k}\Spec k=\A^1$.
There are no extra isomorphism choices, because scheme-valued objects
have only identity arrows.
\item For the atlas $\Spec k\to B\mu_2$, the fibre product
$\Spec k\times_{B\mu_2}\Spec k$ is $\mu_2$, consisting of two points.
The two points are the two isomorphisms between the trivial torsor and itself.
\end{enumerate}}

\examplegloss{flat-and-faithfully-flat}{Flat and faithfully flat}{
\begin{enumerate}[label=(\arabic*),leftmargin=*,itemsep=2pt,topsep=3pt]
\item The map $\A^1_z\to\A^1_u$, $u=z^2$, is faithfully flat.
Indeed, $k[z]$ is a free $k[u]$-module with basis $1,z$, and the
map is surjective, including over $u=0$.
\item The open immersion $\Gm\hookrightarrow\A^1$ is flat because
its ring map is localisation $k[z]\to k[z,z^{-1}]$. It is not
faithfully flat: its image omits the origin.
\end{enumerate}}

\examplegloss{fppf-topology}{Fppf topology}{
\begin{enumerate}[label=(\arabic*),leftmargin=*,itemsep=2pt,topsep=3pt]
\item The single map $\A^1_z\to\A^1_u$, $u=z^2$, is an fppf cover:
it is finite, flat, and surjective. It is not an étale cover, since
it is ramified at zero.
\item Restricting to $\Gm$ gives $s\mapsto s^2$, which is both
étale and fppf. Over this cover the unit $t$ on the base has the
chosen square root $s$.
\end{enumerate}}

\examplegloss{free-action}{Free action}{
\begin{enumerate}[label=(\arabic*),leftmargin=*,itemsep=2pt,topsep=3pt]
\item Multiplication of $\Gm$ on itself is free: $tx=x$ with
$x\ne0$ forces $t=1$. The quotient stack $[\Gm/\Gm]$ is
the ordinary one-point scheme.
\item The sign action of $\mu_2$ on $\Gm$ is free, since
$-z=z$ is impossible for $z\ne0$ in characteristic zero.
The same action on $\A^1$ is not free: zero has stabiliser $\mu_2$.
\end{enumerate}}

\examplegloss{functor-of-points}{Functor of points}{
\begin{enumerate}[label=(\arabic*),leftmargin=*,itemsep=2pt,topsep=3pt]
\item A morphism $T\to\A^1$ is exactly a global regular function
on $T$. Thus $\A^1(T)=\Gamma(T,\OO_T)$; pulling back a point
means pulling back that function.
\item For $T=\Spec k[\varepsilon]/(\varepsilon^2)$, the functions
$0$ and $\varepsilon$ give different $T$-points of $\A^1$.
Both restrict to the origin over $\Spec k$, so ordinary geometric
points alone do not distinguish them.
\end{enumerate}}

\examplegloss{geometric-fibre-weight}{Geometric fibre weight}{
\begin{enumerate}[label=(\arabic*),leftmargin=*,itemsep=2pt,topsep=3pt]
\item On a line over a fixed point, the action $t\cdot v=t^2v$
has geometric fibre weight $2$. In the notes' section convention,
a nonzero constant section of this line has weight $-2$.
\item The sign line on $B\mu_2$ has action $v\mapsto-v$.
Its fibre character is $1$ modulo $2$; its tensor square has
trivial character, although the original line does not.
\end{enumerate}}

\examplegloss{geometric-point}{Geometric point}{
\begin{enumerate}[label=(\arabic*),leftmargin=*,itemsep=2pt,topsep=3pt]
\item The assignment $z\mapsto3$ defines a geometric point
$\Spec k\to\A^1_z$. Its only automorphism as a point of
the scheme is the identity.
\item The origin defines a geometric object of $[\A^1/\mu_2]$
for the sign action. It has two automorphisms, $1$ and $-1$.
Its image in the coarse line is an ordinary point without these arrows.
\end{enumerate}}

\examplegloss{git-quotient}{GIT quotient}{
\begin{enumerate}[label=(\arabic*),leftmargin=*,itemsep=2pt,topsep=3pt]
\item For the weight-$(1,1)$ action on $\A^2$, with positive
linearisation, the semistable locus is $\A^2\setminus0$ and
the quotient is $\Proj k[x,y]=\PP^1$. Here both variables
have tensor degree one.
\item For weights $(2,1)$ the quotient stack is $\PP(2,1)$,
while the GIT quotient is $\Proj k[A,B]$ with degrees $(2,1)$.
Its second Veronese is $k[A,B^2]$, giving coarse space $\PP^1$.
\end{enumerate}}

\examplegloss{good-moduli-space}{Good moduli space}{
\begin{enumerate}[label=(\arabic*),leftmargin=*,itemsep=2pt,topsep=3pt]
\item The coarse map $[\A^1/\mu_2]\to\A^1_u$, $u=z^2$,
is also a good moduli space. Functions downstairs are the even
polynomials, and taking invariants is exact by averaging.
\item For scaling, $[\A^1/\Gm]\to\Spec k$ is good: functions
of weight zero are constants, and taking weight-zero parts is exact.
It identifies zero and the nonzero orbit, so it is not coarse.
\end{enumerate}}

\examplegloss{groupoid}{Groupoid}{
\begin{enumerate}[label=(\arabic*),leftmargin=*,itemsep=2pt,topsep=3pt]
\item A set $\{a,b\}$ becomes a groupoid by giving each object
only its identity arrow. It remembers two objects and no nontrivial
symmetries; this is a discrete groupoid.
\item A one-object groupoid with automorphism group $\mu_2$
has two arrows, identity and sign, with sign composed with itself
equal to identity. It has just one isomorphism class but retains a symmetry.
\end{enumerate}}

\examplegloss{groupoid-presentation}{Groupoid presentation}{
\begin{enumerate}[label=(\arabic*),leftmargin=*,itemsep=2pt,topsep=3pt]
\item Using the identity atlas of $\A^1$, the arrow space is
$\A^1\times_{\A^1}\A^1=\A^1$ and both projections are
identity. There is only one arrow at each object.
\item For $B\mu_2$, take object scheme $\Spec k$ and arrow
scheme $\mu_2$. Both source and target go to the one object,
and arrow composition is multiplication in $\mu_2$.
\end{enumerate}}

\examplegloss{inertia-stack}{Inertia stack}{
\begin{enumerate}[label=(\arabic*),leftmargin=*,itemsep=2pt,topsep=3pt]
\item For the scheme $\A^1$, every object has only the identity
automorphism. Consequently $I\A^1\simeq\A^1$: choosing an
automorphism adds no data.
\item For the sign quotient, the identity automorphism is available
everywhere, while sign fixes only zero. Hence
$I[\A^1/\mu_2]\simeq[\A^1/\mu_2]\sqcup B\mu_2$.
The second component records the extra automorphism at the origin.
\end{enumerate}}

\examplegloss{invariant-function-or-section}{Invariant function or section}{
\begin{enumerate}[label=(\arabic*),leftmargin=*,itemsep=2pt,topsep=3pt]
\item Under $z\mapsto-z$, the function $z^2$ is invariant:
$(-z)^2=z^2$. The function $z$ is not invariant, since it
changes sign.
\item Give the trivial line over $\A^1$ sign action on its fibre.
Then $f(z)=z$ is an equivariant section because $f(-z)=-f(z)$.
It defines a section of the line on the quotient stack despite
its coefficient function being noninvariant.
\end{enumerate}}

\examplegloss{isomorphism-sheaf}{Isomorphism sheaf}{
\begin{enumerate}[label=(\arabic*),leftmargin=*,itemsep=2pt,topsep=3pt]
\item Automorphisms of the trivial line on $T$ are multiplication
by invertible functions. Its isomorphism sheaf with itself is
$\Gm\times T$, rather than the additive line $\A^1\times T$:
zero multiplication is not invertible.
\item The trivial $\mu_2$-torsor over $T$ has automorphism sheaf
$\mu_2\times T$. It is the pullback of the diagonal of
$B\mu_2$ along the corresponding pair of identical families.
\end{enumerate}}

\examplegloss{line-bundle}{Line bundle}{
\begin{enumerate}[label=(\arabic*),leftmargin=*,itemsep=2pt,topsep=3pt]
\item The trivial line $\OO_{\A^1}$ has sections $k[z]$ and
one-dimensional fibres. Its nowhere-zero section $1$ gives a
global trivialisation.
\item The sign character defines a line on $B\mu_2$. Its pullback
to the atlas $\Spec k$ is an ordinary one-dimensional vector space,
but the retained sign action makes it a nontrivial stack line.
Its tensor square is trivial.
\end{enumerate}}

\examplegloss{linearisation}{Linearisation}{
\begin{enumerate}[label=(\arabic*),leftmargin=*,itemsep=2pt,topsep=3pt]
\item On $\A^1$ with $t\cdot z=tz$, lift the action to a trivial
line by $(z,v)\mapsto(tz,v)$. An invariant polynomial section
satisfies $f(tz)=f(z)$, so it is constant.
\item On the same underlying line use $(z,v)\mapsto(tz,tv)$.
Invariance now requires $f(tz)=t f(z)$, giving the multiples
of $z$. The bundle is unchanged, but its linearisation changes the sections.
\end{enumerate}}

\examplegloss{linearly-reductive-group}{Linearly reductive group}{
\begin{enumerate}[label=(\arabic*),leftmargin=*,itemsep=2pt,topsep=3pt]
\item For a $\mu_2$-representation with sign operator $\sigma$,
the map $v\mapsto(v+\sigma v)/2$ projects onto the invariant
vectors. Applying this averaging to a lift of an invariant vector
proves exactness of invariants.
\item A rational $\Gm$-representation splits into weight spaces
$V=\bigoplus_jV_j$. Its invariant space is the direct summand
$V_0$; an equivariant surjection is surjective on this summand.
\end{enumerate}}

\examplegloss{locally-free-sheaf}{Locally free sheaf}{
\begin{enumerate}[label=(\arabic*),leftmargin=*,itemsep=2pt,topsep=3pt]
\item The sheaf $\OO_{\A^1}^{\oplus2}$ is locally free of rank
two; the two constant basis sections trivialise it everywhere.
Its fibre at every point is $k^2$.
\item The line $\OO_{\PP^1}(1)$ is locally free of rank one:
each standard affine chart has a local generator, and the generators
are related by an invertible transition function. A global trivialisation
is not required for local freeness.
\end{enumerate}}

\examplegloss{normal-line}{Normal line}{
\begin{enumerate}[label=(\arabic*),leftmargin=*,itemsep=2pt,topsep=3pt]
\item For the origin in $\A^1_z$, the conormal is $(z)/(z^2)$
and its dual is the normal line. It is an ordinary one-dimensional
vector space over $k$.
\item For $B\mu_2\subset[\A^1_z/\mu_2]$, the normal vector
changes sign. This normal line has sign character, whereas the
pullback of the coarse origin's ordinary normal line has trivial character.
\end{enumerate}}

\examplegloss{orbit}{Orbit}{
\begin{enumerate}[label=(\arabic*),leftmargin=*,itemsep=2pt,topsep=3pt]
\item Under sign on $\A^1$, the orbit of $1$ is $\{1,-1\}$,
while the orbit of zero is $\{0\}$. The first has trivial
stabiliser and the second has stabiliser $\mu_2$.
\item Under scaling by $\Gm$, all nonzero points form one orbit.
Its closure is the whole affine line and contains the separate
zero orbit. The invariant affine quotient identifies the two orbits.
\end{enumerate}}

\examplegloss{polarisation}{Polarisation}{
\begin{enumerate}[label=(\arabic*),leftmargin=*,itemsep=2pt,topsep=3pt]
\item On $\PP^1$, both $\OO(1)$ and $\OO(2)$ are ample
polarisations. Their spaces of sections have dimensions two and
three, respectively, although the underlying curve is the same.
\item On the lower-cut stack $\PP(2,1)$, the level-one polarisation
is $\OO(1)$ with sign character at the boundary. It does not
descend; its square $\OO(2)$ is the pullback of coarse $\OO_{\PP^1}(1)$.
\end{enumerate}}

\examplegloss{projective-degree}{Projective degree}{
\begin{enumerate}[label=(\arabic*),leftmargin=*,itemsep=2pt,topsep=3pt]
\item For $M=\A^1_x$ and $L=\OO_M$, write the section ring
as $k[x,h]$, where $\deg_m x=0$ and $\deg_m h=1$.
The function $x$ is a degree-zero element even though it is nonconstant.
\item If $x$ has group weight two and $h$ has weight zero,
then $xh$ has tensor degree one and group weight two;
$h^2$ has tensor degree two and group weight zero.
\end{enumerate}}

\examplegloss{pullback}{Pullback}{
\begin{enumerate}[label=(\arabic*),leftmargin=*,itemsep=2pt,topsep=3pt]
\item For any map $f:T\to\A^1$, the pullback
$f^*\OO_{\A^1}$ is $\OO_T$. The constant generator $1$
pulls back to the constant generator $1$.
\item For the coarse map $\pi:\PP(2,1)\to\PP^1$,
$\pi^*\OO_{\PP^1}(1)=\OO_{\PP(2,1)}(2)$.
Its boundary stabiliser acts trivially, because the fibre character
is $2=0$ modulo $2$.
\end{enumerate}}

\examplegloss{pushforward}{Pushforward}{
\begin{enumerate}[label=(\arabic*),leftmargin=*,itemsep=2pt,topsep=3pt]
\item For $\pi:[\A^1_z/\mu_2]\to\A^1_u$, $u=z^2$,
the structure-sheaf pushforward corresponds to the invariant module
$k[z]^{\mu_2}=k[u]$.
\item The sign line has pushforward module $z k[u]$, the odd
polynomials. This module is free of rank one over $k[u]$,
but the original stack line does not descend: its origin fibre
still has sign character.
\end{enumerate}}

\examplegloss{quasi-coherent-sheaf}{Quasi-coherent sheaf}{
\begin{enumerate}[label=(\arabic*),leftmargin=*,itemsep=2pt,topsep=3pt]
\item The module $k[z]$ gives the quasi-coherent sheaf
$\OO_{\A^1}$. On the open set $z\ne0$, its module of sections
is the localisation $k[z,z^{-1}]$.
\item The module $k[z]/(z)$ gives a quasi-coherent sheaf supported
at zero. Its fibre has dimension one at zero and zero away
from zero, so it is not a locally free sheaf on $\A^1$.
\end{enumerate}}

\examplegloss{quotient-stack}{Quotient stack}{
\begin{enumerate}[label=(\arabic*),leftmargin=*,itemsep=2pt,topsep=3pt]
\item For multiplication on $\Gm$, the quotient
$[\Gm/\Gm]\simeq\Spec k$. The torsor description recovers
an ordinary quotient because the action is free and transitive.
\item For the trivial action of $\mu_2$ on a point, the quotient
is $B\mu_2$. Over a test scheme it remembers a $\mu_2$-torsor;
over $k$ its one isomorphism class retains two automorphisms.
\end{enumerate}}

\examplegloss{regular-cutting-level}{Regular cutting level}{
\begin{enumerate}[label=(\arabic*),leftmargin=*,itemsep=2pt,topsep=3pt]
\item On $M=\A^1$ with point weight one and trivially
linearised $L=\OO_M$, the only fixed-component section weight
is zero. Thus $d=1$ is regular, and the lower cut is the
ordinary projective line.
\item Replace the point weight by two. The same $d=1$ is
still regular, but the lower cut is $\PP(2,1)$ with boundary
stabiliser $\mu_2$. Regularity allows this finite stabiliser.
\end{enumerate}}

\examplegloss{representable-morphism}{Representable morphism}{
\begin{enumerate}[label=(\arabic*),leftmargin=*,itemsep=2pt,topsep=3pt]
\item The atlas $\A^1\to[\A^1/\mu_2]$ is representable.
After a base change $T\to[\A^1/\mu_2]$, it becomes the
torsor $P\to T$, which is a finite étale morphism of schemes.
\item In contrast, $B\mu_2\to\Spec k$ is not representable:
base changing by $\Spec k$ leaves $B\mu_2$, whose nontrivial
automorphisms prevent it from being an algebraic space.
\end{enumerate}}
